\documentclass[a4paper,11pt]{article}

% ---- Encoding & fonts -------------------------------------------------------
\usepackage[T1]{fontenc}
\usepackage[utf8]{inputenc}
\usepackage{lmodern}
\usepackage{microtype}              % protrusion/kerning -> fewer overfull boxes

% ---- Geometry (~1in margins) ------------------------------------------------
\usepackage[a4paper,margin=1in,headheight=15pt,headsep=14pt]{geometry}

% ---- Math -------------------------------------------------------------------
\usepackage{amsmath,amssymb}

% ---- Graphics, tables, colour ----------------------------------------------
\usepackage{graphicx}
\usepackage{booktabs}
\usepackage[table,svgnames]{xcolor}

% ---- Callout boxes (breakable pill-tab tcolorboxes) -------------------------
\usepackage[most]{tcolorbox}
\tcbuselibrary{skins,breakable}

% ---- Callout glyphs (robust + cross-platform) -------------------------------
\usepackage{pifont}
\IfFileExists{bbding.sty}{%
  \usepackage{bbding}%
  \newcommand{\glyphHand}{\Pointinghand}%
  \newcommand{\glyphPencil}{\Pencil}%
}{%
  \newcommand{\glyphHand}{\ding{43}}%
  \newcommand{\glyphPencil}{\ding{46}}%
}

% ---- Lists & spacing --------------------------------------------------------
\usepackage{enumitem}
\setlist{leftmargin=1.5em,topsep=4pt,itemsep=3pt}

% ---- Dedicated "Step 1. / Step 2. ..." list for WORKED EXAMPLES -------------
\newlist{steps}{enumerate}{1}
\setlist[steps]{label=\textbf{Step \arabic*.},
  leftmargin=4.4em, labelwidth=3.8em, labelsep=0.6em, labelindent=0pt,
  align=left, topsep=5pt, itemsep=6pt}
\usepackage{parskip}                % blank-line paragraph spacing, no indent

% ---- Headers / footers ------------------------------------------------------
\usepackage{fancyhdr}

% ---- Section styling --------------------------------------------------------
\usepackage{titlesec}

% ---- Links (hidden) + URL line-breaking so long URLs never overflow ---------
\usepackage[hidelinks]{hyperref}
\usepackage{xurl}
\hypersetup{breaklinks=true}

% =============================================================================
%  COLOURS — navy chrome + the FIVE taxonomy ramps (frame + light fill)
% =============================================================================
\definecolor{navy}{HTML}{21355E}        % banner + headings
\definecolor{navyrule}{HTML}{21355E}
\definecolor{inktext}{HTML}{1A202C}     % body ink
\definecolor{mutedtext}{HTML}{6B7280}   % header / meta grey

% Intuition — blue
\definecolor{intuFrame}{HTML}{2C5AA0}   \definecolor{intuFill}{HTML}{F4F7FC}  % ~ blue!4
% Everyday — amber
\definecolor{everyFrame}{HTML}{8A5A1E}  \definecolor{everyFill}{HTML}{FBF1E3}  % ~ orange!7
% Worked — green
\definecolor{workFrame}{HTML}{2E7D52}   \definecolor{workFill}{HTML}{F1F8F3}  % ~ green!5
% Watch out — red
\definecolor{watchFrame}{HTML}{B23A48}  \definecolor{watchFill}{HTML}{FBF1F2}  % ~ red!4
% Key takeaway — purple/violet
\definecolor{keyFrame}{HTML}{6A4C93}    \definecolor{keyFill}{HTML}{F4F0F9}    % ~ violet!6
% Wait, what? — teal (the honest surprise)
\definecolor{waitFrame}{HTML}{0F766E}   \definecolor{waitFill}{HTML}{EEF7F5}   % ~ teal!5

\color{inktext}
\hypersetup{linkcolor=navy,urlcolor=navy,citecolor=navy}

% =============================================================================
%  RUNNING HEADER / FOOTER  (fancyhdr)
% =============================================================================
\newcommand{\CourseShort}{DNN AIML ZG511}
\newcommand{\SessionNo}{9}
\newcommand{\LectureTitle}{CNN Architectures, Transfer Learning, RNNs, GRU \& LSTM}

\pagestyle{fancy}
\fancyhf{}
\fancyhead[L]{\footnotesize\color{mutedtext}\textsf{\CourseShort}}
\fancyhead[R]{\footnotesize\color{mutedtext}\textsf{Session \SessionNo\ \textperiodcentered\ \LectureTitle}}
\fancyfoot[L]{\footnotesize\color{mutedtext}\textsf{WILP, BITS Pilani}}
\fancyfoot[C]{\footnotesize\color{mutedtext}\thepage}
\renewcommand{\headrulewidth}{0.4pt}
\renewcommand{\footrulewidth}{0pt}
\renewcommand{\headrule}{\hbox to\headwidth{\color{navyrule!55}\leaders\hrule height \headrulewidth\hfill}}

% =============================================================================
%  SECTION HEADINGS (titlesec)
% =============================================================================
\titleformat{\section}
  {\sffamily\Large\bfseries\color{navy}}
  {\thesection}
  {0.8em}
  {}
  [{\vspace{2pt}\color{navy!70}\titlerule[0.8pt]}]
\titleformat{\subsection}
  {\sffamily\large\bfseries\color{navy!92}}
  {\thesubsection}{0.7em}{}
\titlespacing*{\section}{0pt}{16pt}{8pt}
\titlespacing*{\subsection}{0pt}{12pt}{4pt}

\newif\ifmlkfirstsec \mlkfirstsectrue
\newcommand{\sectionbreak}{\ifmlkfirstsec\mlkfirstsecfalse\else\clearpage\fi}

\newcommand{\secsub}[1]{%
  \noindent{\sffamily\bfseries\itshape\color{navy!85}\large #1}\par\medskip}

\newcommand{\flipanswer}[1]{\rotatebox[origin=c]{180}{\parbox{0.92\linewidth}{\small #1}}}

% =============================================================================
%  PILL-TAB CALLOUTS
% =============================================================================
\tcbset{
  housecallout/.style 2 args={
    enhanced, breakable,
    colback=#2, colframe=#1,
    boxrule=0.6pt, arc=2.4pt, outer arc=2.4pt,
    left=10pt, right=10pt, top=9pt, bottom=9pt,
    before skip=11pt, after skip=11pt,
    fontupper=\normalsize,
    coltitle=white,
    fonttitle=\bfseries\sffamily\small,
    attach boxed title to top left={xshift=6mm,yshift=-3mm},
    boxed title style={
      colback=#1, colframe=#1,
      rounded corners, arc=3pt,
      boxrule=0pt, left=6pt, right=6pt, top=2pt, bottom=2pt
    }
  }
}

\newtcolorbox{intuition}{%
  housecallout={intuFrame}{intuFill},
  title={\raisebox{-0.05ex}{\glyphHand}\,\ The intuition}}

\newtcolorbox{everyday}{%
  housecallout={everyFrame}{everyFill},
  title={\raisebox{-0.05ex}{$\bigstar$}\,\ Everyday picture}}

\newtcolorbox{worked}[1]{%
  housecallout={workFrame}{workFill},
  title={\raisebox{-0.05ex}{\glyphPencil}\,\ Worked example: #1}}

\newtcolorbox{watchout}{%
  housecallout={watchFrame}{watchFill},
  title={\raisebox{-0.05ex}{\ding{55}}\,\ Watch out}}

\newtcolorbox{keytake}{%
  housecallout={keyFrame}{keyFill},
  title={\raisebox{-0.05ex}{\ding{51}}\,\ Key takeaway}}

\newtcolorbox{waitwhat}{%
  housecallout={waitFrame}{waitFill},
  title={\raisebox{-0.05ex}{\ding{93}}\,\ Wait --- really?}}

% =============================================================================
%  BITS BRAND
% =============================================================================
\newif\ifbitslogo \bitslogofalse
\IfFileExists{assets/bits-logo.png}%
  {\bitslogotrue \def\bitslogopath{assets/bits-logo.png}}%
  {\IfFileExists{../../templates/assets/bits-logo.png}%
     {\bitslogotrue \def\bitslogopath{../../templates/assets/bits-logo.png}}%
     {}}
\newcommand{\bitswordmark}{{\sffamily\bfseries\color{white}\fontsize{12.5}{15}\selectfont WILP, BITS Pilani}}
\newcommand{\bitslogochip}{\ifbitslogo\colorbox{white}{\,\raisebox{-0.2ex}{\includegraphics[height=26pt]{\bitslogopath}}\,}\fi}

% =============================================================================
%  TITLE BANNER
% =============================================================================
\providecommand{\addfontfeatureslike}[1]{#1}
\newcommand{\titlebanner}[4]{%
  \begin{tcolorbox}[
    enhanced, breakable=false,
    colback=navy, colframe=navy,
    boxrule=0pt, arc=3pt, outer arc=3pt,
    left=16pt, right=16pt, top=16pt, bottom=16pt,
    before skip=2pt, after skip=14pt,
    width=\linewidth ]
    \noindent
    \begin{minipage}[c]{0.72\linewidth}\bitswordmark\end{minipage}%
    \hfill
    \begin{minipage}[c]{0.26\linewidth}\raggedleft\bitslogochip\end{minipage}\par
    \vspace{9pt}
    {\sffamily\footnotesize\bfseries\color{white!85}%
      \MakeUppercase{\addfontfeatureslike{#1}}}\par
    \vspace{6pt}
    {\sffamily\bfseries\fontsize{26}{30}\selectfont\color{white} #2\par}
    \vspace{5pt}
    {\sffamily\large\color{white!88} #3\par}
    \vspace{9pt}
    {\color{white!55}\rule{\linewidth}{0.4pt}}\par
    \vspace{6pt}
    {\sffamily\footnotesize\color{white!82} #4\par}
  \end{tcolorbox}%
}

\newcommand{\housefig}[2]{%
  \begin{figure}[htbp]
    \centering
    \includegraphics[width=\linewidth]{#1}%
    \caption{#2}
  \end{figure}%
}

\newcommand{\vocab}[1]{\textbf{\color{navy}#1}}

\begin{document}

% ----------------------------------------------------------------------------
%  PAGE 1 — TITLE BANNER
% ----------------------------------------------------------------------------
\thispagestyle{fancy}

\titlebanner
  {Deep Neural Networks (AIML ZG511) \textperiodcentered\ Companion Reader}% LABEL
  {CNN Patterns, Transfer Learning \& Sequence Models}% BIG TITLE
  {From Spatial Features and ResNets to RNNs, BPTT, GRU, and LSTM Memory Cells}% SUBTITLE
  {Session \SessionNo\ (5th July 2026) \textperiodcentered\ Read alongside the lecture slides \textperiodcentered\ Every slide example worked out in full} % META

\smallskip
\noindent{\small\color{mutedtext}%
Prepared for the Work Integrated Learning Programmes (WILP), BITS Pilani.}
\smallskip

\noindent\textbf{How to use this companion.}
This reader expands Session 9 into a complete study guide. We cover two core themes: advanced \vocab{Convolutional Neural Network (CNN)} design patterns and \vocab{Recurrent Neural Network (RNN)} sequence architectures. Every formula is derived step by step. Every slide example is calculated with real numbers. Use the colour-coded boxes for intuition, everyday analogies, worked math, and common pitfalls.

\medskip

% ============================================================================
%  SECTION 1: OPENING BY DOING
% ============================================================================
\section{Opening by doing: Spatial features vs temporal memory}
\secsub{Try a 2-step sequence calculation by hand before defining recurrent states.}

Consider a small factory line that processes incoming data. An image model looks at a whole grid at once. A sequence model reads items one by one over time. Let us calculate how a simple temporal hidden state updates over 2 steps right now.

Suppose your initial memory state is $h_0 = [0, 0]^T$. At time step $t=1$, an input $x_1 = [1, 0]^T$ arrives. The memory weight matrix is $W_{hh} = \begin{bmatrix} 0.5 & 0.0 \\ 0.0 & 0.5 \end{bmatrix}$ and input weight matrix is $W_{hx} = \begin{bmatrix} 0.8 & 0.0 \\ 0.0 & 0.8 \end{bmatrix}$.

The raw hidden state before activation is $z_1 = W_{hh} h_0 + W_{hx} x_1 = [0, 0]^T + [0.8, 0.0]^T = [0.8, 0.0]^T$. Applying $\tanh$ gives $h_1 = [\tanh(0.8), \tanh(0.0)]^T = [0.664, 0.0]^T$.

At time step $t=2$, the input changes to $x_2 = [0, 1]^T$. The new state incorporates past memory $h_1$:
\begin{equation}
  z_2 = W_{hh} h_1 + W_{hx} x_2 = \begin{bmatrix} 0.5 & 0 \\ 0 & 0.5 \end{bmatrix} \begin{bmatrix} 0.664 \\ 0 \end{bmatrix} + \begin{bmatrix} 0.8 & 0 \\ 0 & 0.8 \end{bmatrix} \begin{bmatrix} 0 \\ 1 \end{bmatrix} = \begin{bmatrix} 0.332 \\ 0.800 \end{bmatrix}.
\end{equation}
Then $h_2 = [\tanh(0.332), \tanh(0.800)]^T = [0.321, 0.664]^T$.

Notice what just happened: $h_2$ carries information from \emph{both} $x_1$ and $x_2$. That single matrix multiplication $W_{hh} h_{t-1}$ is the magic ingredient that gives neural networks temporal memory!

% ============================================================================
%  SECTION 2: CNN ARCHITECTURAL PATTERNS
% ============================================================================
\section{CNN design patterns and architectural decisions}
\secsub{Five reusable design patterns turn arbitrary layers into scalable, high-performance visual networks.}

A typical Convolutional Neural Network alternates between convolution, non-linear activation, and spatial pooling. Early layers capture low-level features such as edges and corners. Middle layers combine edges into textures and shapes. Later layers form high-level concepts like object parts or entire faces. To build deep networks reliably, five key design patterns are used in modern deep learning:

\begin{enumerate}[leftmargin=1.5em]
  \item \vocab{Pattern 1: Spatial Reduction + Channel Expansion.} As spatial dimensions decrease ($32 \times 32 \to 14 \times 14 \to 5 \times 5$), feature channels increase ($3 \to 16 \to 64 \to 256$). This preserves computational budget while increasing representational capacity.
  \item \vocab{Pattern 2: Repetitive Building Blocks.} Replace custom layer arrangements with repeated modular blocks (e.g., VGG blocks or ResNet blocks) for easier scaling.
  \item \vocab{Pattern 3: Skip / Residual Connections.} Pass inputs directly across layers ($y = F(x) + x$). If a layer learns nothing, it defaults to the identity mapping. This solves vanishing gradients in networks deeper than 100 layers.
  \item \vocab{Pattern 4: Bottleneck Designs.} Use $1 \times 1$ convolutions to squash channel depth before expensive $3 \times 3$ convolutions, then expand depth back. This reduces parameters by up to 8x.
  \item \vocab{Pattern 5: Multi-Scale Feature Extraction.} Use parallel filters of different sizes ($1 \times 1, 3 \times 3, 5 \times 5$) inside Inception modules to capture features at multiple spatial scales simultaneously.
\end{enumerate}

\begin{intuition}
A residual block allows gradients to flow down an express highway during backpropagation. Instead of forcing layers to re-learn the full signal, the layer only learns the \emph{change} or residual $F(x)$.
\end{intuition}

\begin{everyday}
Think of an express train line alongside local tracks. The local tracks ($F(x)$) perform detailed station-by-station work. The express track ($x$) bypasses stations entirely, letting passenger signals travel directly to the destination without getting lost.
\end{everyday}

\housefig{figures/cnn_patterns.png}{CNN Design Patterns: Spatial reduction, channel expansion, bottleneck parameter reduction, and residual skip connections.}

\begin{worked}{Bottleneck vs Standard Convolution Parameter Reduction}
Compare a standard convolution layer against a bottleneck block operating on a tensor with 256 channels.
\begin{steps}
  \item \textbf{Standard Conv 3x3.} Input: 256 channels, Output: 256 channels.
        $$\text{Parameters} = 3 \times 3 \times 256 \times 256 = 589,824 \approx 590\text{K}.$$
  \item \textbf{Bottleneck Block.}
        \begin{itemize}
          \item Step A ($1 \times 1$ conv to 64 channels): $1 \times 1 \times 256 \times 64 = 16,384$.
          \item Step B ($3 \times 3$ conv on 64 channels): $3 \times 3 \times 64 \times 64 = 36,864$.
          \item Step C ($1 \times 1$ conv to 256 channels): $1 \times 1 \times 64 \times 256 = 16,384$.
        \end{itemize}
        $$\text{Total Parameters} = 16,384 + 36,864 + 16,384 = 69,632 \approx 70\text{K}.$$
  \item \textbf{Reduction ratio.} $589,824 / 69,632 \approx 8.47\times$ parameter reduction!
\end{steps}
\end{worked}

\begin{watchout}
Depth alone is not enough. Adding layers to a plain feedforward network without skip connections or batch normalization causes training accuracy to degrade rapidly due to vanishing gradients.
\end{watchout}

\begin{keytake}
Effective CNN architectures combine spatial reduction with channel expansion, standard modular blocks, residual connections, and bottleneck layers.
\end{keytake}

% ============================================================================
%  SECTION 3: CLASSIC CNN ARCHITECTURES & TRANSFER LEARNING
% ============================================================================
\section{Classic CNN architectures and transfer learning}
\secsub{AlexNet, VGG, and GoogLeNet laid the groundwork for modern computer vision and transfer learning.}

The evolution of ImageNet models highlights key architectural shifts:
\begin{itemize}
  \item \vocab{AlexNet (2012):} 8 layers, ~60M parameters. Introduced ReLU, Local Response Normalization (LRN), Dropout (0.5), heavy data augmentation, and GPU training.
  \item \vocab{VGG-16 (2014):} 16 layers, ~138M parameters. Proved that stacking simple $3 \times 3$ convolutions with $2 \times 2$ max pooling yields deeper, richer features.
  \item \vocab{GoogLeNet / Inception (2014):} 22 layers, ~5M parameters (12x fewer parameters than AlexNet). Used Inception modules and replaced heavy Fully Connected (FC) layers with Global Average Pooling (GAP).
\end{itemize}

\begin{table}[htbp]
  \centering
  \caption{Comparison of Classic Vision Architectures}
  \resizebox{\linewidth}{!}{%
  \begin{tabular}{lcccc}
    \toprule
    \textbf{Architecture} & \textbf{Depth (Layers)} & \textbf{Parameters} & \textbf{Key Innovation} & \textbf{Classifier} \\
    \midrule
    AlexNet & 8 & ~60M & ReLU, Dropout, LRN & 3 FC Layers \\
    VGG-16 & 16 & ~138M & Small $3 \times 3$ filters stacked & 3 FC Layers \\
    GoogLeNet & 22 & ~5M & Inception parallel blocks & Global Avg Pool \\
    ResNet-50 & 50 & ~25.6M & Residual skip connections & Global Avg Pool \\
    \bottomrule
  \end{tabular}%
  }
\end{table}

\noindent\textbf{Transfer Learning Strategy.}
Instead of training from scratch, pre-trained weights on large datasets (like ImageNet) can be transferred to new target tasks:
\begin{enumerate}[leftmargin=1.5em]
  \item \textbf{Small target dataset + high similarity:} Freeze pre-trained backbone; train linear classifier on top features.
  \item \textbf{Large target dataset + high similarity:} Fine-tune all layers with a small learning rate.
  \item \textbf{Small target dataset + low similarity:} Freeze early layers; train higher layers and linear head.
  \item \textbf{Large target dataset + low similarity:} Train whole network from scratch or fine-tune completely.
\end{enumerate}

\begin{keytake}
Transfer learning uses low-level feature representations learned from massive datasets, reducing training time and data requirements for specialized tasks.
\end{keytake}

% ============================================================================
%  SECTION 4: SEQUENCE MODELING & AUTOREGRESSIVE MODELS
% ============================================================================
\section{Sequence modeling: Autoregressive vs latent variable models}
\secsub{Standard feedforward networks fail on sequence data because inputs vary in length and lack temporal memory.}

Standard feedforward networks assume inputs are independent and identically distributed (i.i.d.) with fixed spatial dimensions. Sequence tasks (e.g., text, speech, stock market, temperature tracking) break these assumptions:
\begin{itemize}
  \item Inputs and outputs have variable lengths.
  \item Order matters: "Not good, bad" vs "Good, not bad".
  \item Temporal parameters must be shared across time steps.
\end{itemize}

To model sequence data $x_1, x_2, \dots, x_T$, two primary approaches exist:

\begin{enumerate}[leftmargin=1.5em]
  \item \vocab{Autoregressive Models:} Predict current step $x_t$ using a fixed window of past observations:
        $$x_t = P(x_t \mid x_{t-1}, x_{t-2}, \dots, x_{t-\tau}).$$
        Under a first-order \vocab{Markov Assumption}, $P(x_t \mid x_{t-1}, \dots, x_1) = P(x_t \mid x_{t-1})$.
  \item \vocab{Latent Variable Models (RNNs):} Maintain an internal hidden state $h_t$ that summarizes the entire history up to time $t$:
        $$h_t = g(h_{t-1}, x_t), \quad \hat{y}_t = f(h_t).$$
\end{enumerate}

\housefig{figures/autoregressive_vs_latent.png}{Autoregressive prediction using past observations vs Latent variable models tracking hidden state history.}

\begin{worked}{Temperature Forecasting — Autoregressive vs Latent Variable Model}
Consider a temperature sequence: $20^\circ\text{C}, 22^\circ\text{C}, 25^\circ\text{C}, 24^\circ\text{C}, 21^\circ\text{C}$. Predict tomorrow's temperature.
\begin{steps}
  \item \textbf{Autoregressive First-Order Model.} Formula: $T_{t+1} = \alpha T_t + \beta$.
        Given learned weights $\alpha = 0.8, \beta = 4.0$, and today's temperature $T_5 = 21^\circ\text{C}$:
        $$T_{\text{tomorrow}} = 0.8 \times 21 + 4.0 = 16.8 + 4.0 = 20.8^\circ\text{C}.$$
  \item \textbf{Latent Variable Recurrent Model.} Formula: $h_t = \tanh(W_{hh} h_{t-1} + W_{hx} T_t + b_h)$.
        The hidden state $h_5$ integrates the entire sequence trend ($20 \to 22 \to 25 \to 24 \to 21$) rather than relying solely on today's single reading.
\end{steps}
\end{worked}

\begin{keytake}
Latent variable models compress long sequence histories into a fixed-size hidden state vector $h_t$, avoiding the need for fixed context windows.
\end{keytake}

% ============================================================================
%  SECTION 5: VANILLA RECURRENT NEURAL NETWORKS (RNNs)
% ============================================================================
\section{Vanilla Recurrent Neural Networks (RNNs)}
\secsub{Unrolling an RNN in time reveals a deep network with parameter sharing across all time steps.}

A vanilla Recurrent Neural Network processes input sequence $x_1, x_2, \dots, x_T$ sequentially. At step $t$, it updates hidden state $h_t$ and computes output $o_t$:
\begin{align}
  h_t &= \phi(W_{hh} h_{t-1} + W_{hx} x_t + b_h), \\
  o_t &= W_{ho} h_t + b_o, \\
  \hat{y}_t &= \text{softmax}(o_t) \quad (\text{or linear for regression}).
\end{align}

Here $\phi(\cdot)$ is typically the hyperbolic tangent function ($\tanh$). Parameters $W_{hh}, W_{hx}, W_{ho}, b_h, b_o$ are shared identically across every time step.

\housefig{figures/rnn_unrolled.png}{Unrolled RNN computation graph showing shared weight matrices $W_{hx}, W_{hh}, W_{ho}$ across time.}

\begin{worked}{RNN Parameter Counting}
Calculate total parameters for a vanilla RNN with input dimension $d=100$, hidden dimension $h=256$, and vocabulary/output size $V=10,000$.
\begin{steps}
  \item \textbf{Hidden layer weight matrices and bias:}
        \begin{itemize}
          \item $W_{hx} \in \mathbb{R}^{h \times d}$: $256 \times 100 = 25,600$.
          \item $W_{hh} \in \mathbb{R}^{h \times h}$: $256 \times 256 = 65,536$.
          \item $b_h \in \mathbb{R}^{h}$: $256$.
        \end{itemize}
  \item \textbf{Output layer weight matrix and bias:}
        \begin{itemize}
          \item $W_{ho} \in \mathbb{R}^{V \times h}$: $10,000 \times 256 = 2,560,000$.
          \item $b_o \in \mathbb{R}^{V}$: $10,000$.
        \end{itemize}
  \item \textbf{Total parameters:}
        $$\text{Total} = 25,600 + 65,536 + 256 + 2,560,000 + 10,000 = 2,661,392 \approx 2.66\text{M}.$$
\end{steps}
\end{worked}

\begin{worked}{Step-by-Step RNN Forward Pass (2 Time Steps)}
Given parameters:
$$W_{hx} = \begin{bmatrix} 0.5 & 0.3 \\ 0.2 & 0.4 \\ 0.1 & 0.2 \end{bmatrix}, \quad W_{hh} = \begin{bmatrix} 0.3 & 0.2 & 0.1 \\ 0.1 & 0.4 & 0.2 \\ 0.2 & 0.1 & 0.3 \end{bmatrix}, \quad b_h = \begin{bmatrix} 0.1 \\ 0.1 \\ 0.1 \end{bmatrix},$$
$$W_{ho} = \begin{bmatrix} 0.2 & 0.6 & 0.1 \\ 0.4 & 0.3 & 0.5 \end{bmatrix}, \quad b_o = \begin{bmatrix} 0.1 \\ 0.1 \end{bmatrix}, \quad h_0 = \begin{bmatrix} 0 \\ 0 \\ 0 \end{bmatrix}.$$
Input sequence: $x_1 = [1.0, 0.5]^T, x_2 = [0.8, 1.2]^T$.

\begin{steps}
  \item \textbf{Time Step $t=1$ — Hidden State $h_1$:}
        $$z_1 = W_{hh} h_0 + W_{hx} x_1 + b_h = \begin{bmatrix} 0 \\ 0 \\ 0 \end{bmatrix} + \begin{bmatrix} 0.5(1.0)+0.3(0.5) \\ 0.2(1.0)+0.4(0.5) \\ 0.1(1.0)+0.2(0.5) \end{bmatrix} + \begin{bmatrix} 0.1 \\ 0.1 \\ 0.1 \end{bmatrix} = \begin{bmatrix} 0.75 \\ 0.50 \\ 0.30 \end{bmatrix}.$$
        $$h_1 = \tanh(z_1) = \begin{bmatrix} \tanh(0.75) \\ \tanh(0.50) \\ \tanh(0.30) \end{bmatrix} = \begin{bmatrix} 0.635 \\ 0.462 \\ 0.291 \end{bmatrix}.$$
  \item \textbf{Time Step $t=1$ — Output $o_1$:}
        $$o_1 = W_{ho} h_1 + b_o = \begin{bmatrix} 0.2(0.635)+0.6(0.462)+0.1(0.291) + 0.1 \\ 0.4(0.635)+0.3(0.462)+0.5(0.291) + 0.1 \end{bmatrix} = \begin{bmatrix} 0.533 \\ 0.638 \end{bmatrix}.$$
  \item \textbf{Time Step $t=2$ — Hidden State $h_2$:}
        $$W_{hh} h_1 = \begin{bmatrix} 0.3(0.635)+0.2(0.462)+0.1(0.291) \\ 0.1(0.635)+0.4(0.462)+0.2(0.291) \\ 0.2(0.635)+0.1(0.462)+0.3(0.291) \end{bmatrix} = \begin{bmatrix} 0.312 \\ 0.306 \\ 0.261 \end{bmatrix}.$$
        $$W_{hx} x_2 = \begin{bmatrix} 0.5(0.8)+0.3(1.2) \\ 0.2(0.8)+0.4(1.2) \\ 0.1(0.8)+0.2(1.2) \end{bmatrix} = \begin{bmatrix} 0.760 \\ 0.640 \\ 0.320 \end{bmatrix}.$$
        $$z_2 = \begin{bmatrix} 0.312 \\ 0.306 \\ 0.261 \end{bmatrix} + \begin{bmatrix} 0.760 \\ 0.640 \\ 0.320 \end{bmatrix} + \begin{bmatrix} 0.1 \\ 0.1 \\ 0.1 \end{bmatrix} = \begin{bmatrix} 1.172 \\ 1.046 \\ 0.681 \end{bmatrix}.$$
        $$h_2 = \tanh(z_2) = \begin{bmatrix} \tanh(1.172) \\ \tanh(1.046) \\ \tanh(0.681) \end{bmatrix} = \begin{bmatrix} 0.825 \\ 0.780 \\ 0.592 \end{bmatrix}.$$
  \item \textbf{Time Step $t=2$ — Output $o_2$:}
        $$o_2 = W_{ho} h_2 + b_o = \begin{bmatrix} 0.2(0.825)+0.6(0.780)+0.1(0.592) + 0.1 \\ 0.4(0.825)+0.3(0.780)+0.5(0.592) + 0.1 \end{bmatrix} = \begin{bmatrix} 0.792 \\ 0.960 \end{bmatrix}.$$
\end{steps}
\end{worked}

\begin{keytake}
A vanilla RNN uses shared weight matrices across time steps to update hidden state $h_t = \tanh(W_{hh} h_{t-1} + W_{hx} x_t + b_h)$.
\end{keytake}

% ============================================================================
%  SECTION 6: BPTT & VANISHING GRADIENTS
% ============================================================================
\section{Backpropagation Through Time (BPTT) and vanishing gradients}
\secsub{Unrolling an RNN allows standard backpropagation, but long sequences cause gradients to explode or vanish exponentially.}

To train an RNN, loss is calculated across sequence time steps: $L = \sum_{t=1}^T L_t$. \vocab{Backpropagation Through Time (BPTT)} calculates gradients with respect to parameters by propagating backwards through time.

The gradient of loss $L$ with respect to hidden state $h_k$ at step $k < T$ involves a chain product of Jacobian matrices:
\begin{equation}
  \frac{\partial L_T}{\partial h_k} = \frac{\partial L_T}{\partial h_T} \prod_{j=k+1}^T \frac{\partial h_j}{\partial h_{j-1}} = \frac{\partial L_T}{\partial h_T} \prod_{j=k+1}^T \text{diag}(\phi'(z_j)) W_{hh}^T.
\end{equation}

\housefig{figures/bptt_gradient_flow.png}{Exponential decay of gradient magnitude across steps in Backpropagation Through Time.}

\begin{worked}{Numerical Breakdown of Vanishing Gradients}
Suppose the largest eigenvalue of $W_{hh}$ is $\lambda_{\max} = 0.8$, and the activation derivative is $|\phi'(z)| \approx 0.5$.
\begin{steps}
  \item \textbf{Single Step Multiplier.} $\gamma = 0.8 \times 0.5 = 0.4$.
  \item \textbf{After 5 steps ($k=5$).} Gradient ratio $= (0.4)^5 = 0.01024 \approx 1\%$.
  \item \textbf{After 10 steps ($k=10$).} Gradient ratio $= (0.4)^{10} = 0.0001048 \approx 0.01\%$.
  \item \textbf{After 20 steps ($k=20$).} Gradient ratio $= (0.4)^{20} \approx 1.09 \times 10^{-8}$.
  \item \textbf{Conclusion.} The gradient signal from step 20 vanishes completely before reaching step 1. Early layers cannot update their parameters to learn long-term context.
\end{steps}
\end{worked}

\begin{watchout}
If weight eigenvalues are greater than 1 (e.g., $\lambda_{\max} > 1.5$), gradients explode exponentially ($1.5^{20} \approx 3325$), causing numerical instability. Gradient clipping is used to cap exploding gradients.
\end{watchout}

\begin{keytake}
BPTT backpropagates errors across sequence steps. Repeated matrix multiplication by $W_{hh}$ causes gradients to vanish exponentially for long dependencies.
\end{keytake}

% ============================================================================
%  SECTION 7: BIDIRECTIONAL & DEEP STACKED RNNS
% ============================================================================
\section{Bidirectional and deep (stacked) RNN architectures}
\secsub{Bidirectional RNNs capture past and future context, while deep RNNs stack hidden layers to build hierarchical representations.}

Standard RNNs only use past context $x_1, \dots, x_{t-1}$. In tasks like translation or Named Entity Recognition (NER), future words clarify present meaning (e.g., "I met \textbf{Paris} Hilton").

\begin{enumerate}[leftmargin=1.5em]
  \item \vocab{Bidirectional RNN (BRNN):} Runs two separate hidden layers: a forward layer $\overrightarrow{h}_t$ ($t=1 \to T$) and a backward layer $\overleftarrow{h}_t$ ($t=T \to 1$). Outputs combine both: $h_t = [\overrightarrow{h}_t; \overleftarrow{h}_t]$.
  \item \vocab{Deep (Stacked) RNN:} Stacks multiple recurrent layers vertically. Layer $l$ receives hidden states from layer $l-1$ at step $t$, and hidden states from layer $l$ at step $t-1$:
        $$h_t^{(l)} = \tanh(W_{hh}^{(l)} h_{t-1}^{(l)} + W_{hx}^{(l)} h_t^{(l-1)} + b_h^{(l)}).$$
\end{enumerate}

\begin{worked}{Bidirectional RNN Parameter Counting}
Given input $d=100$, forward hidden $h_f=128$, backward hidden $h_b=128$, output $V=10,000$.
\begin{steps}
  \item \textbf{Forward RNN Parameters:} $128 \times 100 + 128 \times 128 + 128 = 12,800 + 16,384 + 128 = 29,312$.
  \item \textbf{Backward RNN Parameters:} $128 \times 100 + 128 \times 128 + 128 = 29,312$.
  \item \textbf{Output Layer Parameters:} Concatenated hidden size $= 128 + 128 = 256$.
        $$W_{ho} \in \mathbb{R}^{10,000 \times 256} \implies 2,560,000, \quad b_o \in \mathbb{R}^{10,000} \implies 10,000.$$
  \item \textbf{Total Parameters:} $29,312 + 29,312 + 2,570,000 = 2,628,624 \approx 2.63\text{M}$.
\end{steps}
\end{worked}

\begin{keytake}
Bidirectional RNNs process text in both directions to capture full sequence context. Deep RNNs stack recurrent layers to learn abstract hierarchy.
\end{keytake}

% ============================================================================
%  SECTION 8: GATED RECURRENT UNIT (GRU)
% ============================================================================
\section{Gated Recurrent Unit (GRU)}
\secsub{GRU uses reset and update gates to control information flow, creating additive gradient shortcuts.}

To combat vanishing gradients without complex memory structures, the \vocab{Gated Recurrent Unit (GRU)} introduces gating mechanisms that selectively retain or forget information over time:

\begin{align}
  r_t &= \sigma(W_r x_t + U_r h_{t-1} + b_r) \quad \text{(Reset Gate)}, \\
  z_t &= \sigma(W_z x_t + U_z h_{t-1} + b_z) \quad \text{(Update Gate)}, \\
  \tilde{h}_t &= \tanh(W_h x_t + U_h (r_t \odot h_{t-1}) + b_h) \quad \text{(Candidate State)}, \\
  h_t &= z_t \odot h_{t-1} + (1 - z_t) \odot \tilde{h}_t \quad \text{(Final State)}.
\end{align}

\housefig{figures/gru_architecture.png}{GRU Architecture: Reset gate $r_t$, update gate $z_t$, and candidate hidden state $\tilde{h}_t$.}

\begin{worked}{Complete 5-Step GRU Numerical Calculation}
Inputs: $d=2, h=2$. Given $x_1 = [1.0, 0.5]^T, h_0 = [0.6, 0.4]^T$.
Weights:
\begin{align*}
W_r &= \begin{bmatrix} 0.3 & 0.2 \\ 0.1 & 0.4 \end{bmatrix}, U_r = \begin{bmatrix} 0.5 & 0.1 \\ 0.2 & 0.3 \end{bmatrix}, b_r = \begin{bmatrix} 0.1 \\ 0.1 \end{bmatrix}, \\
W_z &= \begin{bmatrix} 0.2 & 0.5 \\ 0.4 & 0.1 \end{bmatrix}, U_z = \begin{bmatrix} 0.1 & 0.4 \\ 0.3 & 0.2 \end{bmatrix}, b_z = \begin{bmatrix} 0.2 \\ 0.2 \end{bmatrix}, \\
W_h &= \begin{bmatrix} 0.5 & 0.3 \\ 0.2 & 0.4 \end{bmatrix}, U_h = \begin{bmatrix} 0.4 & 0.2 \\ 0.1 & 0.3 \end{bmatrix}, b_h = \begin{bmatrix} 0.1 \\ 0.1 \end{bmatrix}.
\end{align*}

\begin{steps}
  \item \textbf{Step 1: Reset Gate $r_1$.}
        \begin{align*}
          W_r x_1 + U_r h_0 + b_r &= \begin{bmatrix} 0.4 \\ 0.3 \end{bmatrix} + \begin{bmatrix} 0.34 \\ 0.24 \end{bmatrix} + \begin{bmatrix} 0.1 \\ 0.1 \end{bmatrix} = \begin{bmatrix} 0.84 \\ 0.64 \end{bmatrix}, \\
          r_1 &= \sigma\begin{bmatrix} 0.84 \\ 0.64 \end{bmatrix} = \begin{bmatrix} 0.6985 \\ 0.6547 \end{bmatrix}.
        \end{align*}
  \item \textbf{Step 2: Update Gate $z_1$.}
        \begin{align*}
          W_z x_1 + U_z h_0 + b_z &= \begin{bmatrix} 0.45 \\ 0.45 \end{bmatrix} + \begin{bmatrix} 0.22 \\ 0.26 \end{bmatrix} + \begin{bmatrix} 0.2 \\ 0.2 \end{bmatrix} = \begin{bmatrix} 0.87 \\ 0.91 \end{bmatrix}, \\
          z_1 &= \sigma\begin{bmatrix} 0.87 \\ 0.91 \end{bmatrix} = \begin{bmatrix} 0.7047 \\ 0.7130 \end{bmatrix}.
        \end{align*}
  \item \textbf{Step 3: Gated Memory $r_1 \odot h_0$.}
        $$r_1 \odot h_0 = \begin{bmatrix} 0.6985 \times 0.6 \\ 0.6547 \times 0.4 \end{bmatrix} = \begin{bmatrix} 0.4191 \\ 0.2619 \end{bmatrix}.$$
  \item \textbf{Step 4: Candidate Hidden State $\tilde{h}_1$.}
        \begin{align*}
          W_h x_1 + U_h (r_1 \odot h_0) + b_h &= \begin{bmatrix} 0.65 \\ 0.40 \end{bmatrix} + \begin{bmatrix} 0.2200 \\ 0.1205 \end{bmatrix} + \begin{bmatrix} 0.1 \\ 0.1 \end{bmatrix} = \begin{bmatrix} 0.9700 \\ 0.6205 \end{bmatrix}, \\
          \tilde{h}_1 &= \tanh\begin{bmatrix} 0.9700 \\ 0.6205 \end{bmatrix} = \begin{bmatrix} 0.7487 \\ 0.5516 \end{bmatrix}.
        \end{align*}
  \item \textbf{Step 5: Final Hidden State Update $h_1$.}
        \begin{align*}
          h_1 &= z_1 \odot h_0 + (1 - z_1) \odot \tilde{h}_1 \\
              &= \begin{bmatrix} 0.7047(0.6) + (0.2953)(0.7487) \\ 0.7130(0.4) + (0.2870)(0.5516) \end{bmatrix} = \begin{bmatrix} 0.6439 \\ 0.4435 \end{bmatrix}.
        \end{align*}
\end{steps}
\end{worked}

\begin{worked}{GRU Parameter Count Formula}
For input $d$ and hidden size $h$, GRU has 3 gate components ($r, z, \tilde{h}$), each with matrices $W \in \mathbb{R}^{h \times d}, U \in \mathbb{R}^{h \times h}, b \in \mathbb{R}^h$.
$$\text{Parameters} = 3 \times (h \cdot d + h^2 + h) + V \cdot h + V.$$
For $d=100, h=256, V=10,000$:
\begin{align*}
\text{Recurrent Params} &= 3 \times (256 \times 100 + 256^2 + 256) \\
&= 3 \times (25,600 + 65,536 + 256) = 274,176 \approx 274\text{K}.
\end{align*}
\end{worked}

\begin{keytake}
GRU uses an update gate $z_t$ to create an additive shortcut $h_t = z_t \odot h_{t-1} + (1-z_t) \odot \tilde{h}_t$, allowing gradients to pass unchanged across long distances when $z_t \approx 1$.
\end{keytake}

% ============================================================================
%  SECTION 9: LONG SHORT-TERM MEMORY (LSTM)
% ============================================================================
\section{Long Short-Term Memory (LSTM) networks}
\secsub{LSTM separates long-term cell state memory $c_t$ from short-term hidden output $h_t$ using three dedicated gates.}

To maintain precise, long-range temporal dependencies over hundreds of steps, the \vocab{Long Short-Term Memory (LSTM)} architecture splits state memory into two distinct vectors:
\begin{enumerate}[leftmargin=1.5em]
  \item \vocab{Cell State ($c_t$):} The long-term memory conveyor belt.
  \item \vocab{Hidden State ($h_t$):} The short-term working output.
\end{enumerate}

LSTM operates via three dedicated gates:
\begin{align}
  f_t &= \sigma(W_f x_t + U_f h_{t-1} + b_f) \quad \text{(Forget Gate — what to discard)}, \\
  i_t &= \sigma(W_i x_t + U_i h_{t-1} + b_i) \quad \text{(Input Gate — what to update)}, \\
  \tilde{c}_t &= \tanh(W_c x_t + U_c h_{t-1} + b_c) \quad \text{(Candidate Cell State)}, \\
  c_t &= f_t \odot c_{t-1} + i_t \odot \tilde{c}_t \quad \text{(Cell State Update)}, \\
  o_t &= \sigma(W_o x_t + U_o h_{t-1} + b_o) \quad \text{(Output Gate — what to expose)}, \\
  h_t &= o_t \odot \tanh(c_t) \quad \text{(Final Hidden Output)}.
\end{align}

\housefig{figures/lstm_architecture.png}{LSTM Architecture: Cell state $c_t$ highway with Forget gate $f_t$, Input gate $i_t$, and Output gate $o_t$.}

\begin{worked}{LSTM Parameter Count Formula}
LSTM contains 4 gate components ($f, i, o, \tilde{c}$), each with weight matrices $W \in \mathbb{R}^{h \times d}, U \in \mathbb{R}^{h \times h}, b \in \mathbb{R}^h$.
$$\text{Parameters} = 4 \times (h \cdot d + h^2 + h) + V \cdot h + V.$$
For $d=100, h=256, V=10,000$:
\begin{align*}
\text{Recurrent Params} &= 4 \times (256 \times 100 + 256^2 + 256) \\
&= 4 \times (25,600 + 65,536 + 256) = 365,568 \approx 366\text{K}.
\end{align*}
\end{worked}

\begin{waitwhat}
In an LSTM cell state update $c_t = f_t \odot c_{t-1} + i_t \odot \tilde{c}_t$, if $f_t = 1$ and $i_t = 0$, the cell state $c_t = c_{t-1}$ remains completely unchanged without \emph{any} decay! Gradients flow backward along $\frac{\partial c_t}{\partial c_{t-1}} = 1.0$ infinitely without vanishing. This constant error carousel is why LSTMs can bridge gaps of over 100 time steps easily.
\end{waitwhat}

\begin{keytake}
LSTM isolates long-term memory inside cell state $c_t$. It uses forget, input, and output gates to achieve constant error flow across long time horizons.
\end{keytake}

% ============================================================================
%  SECTION 10: SYNTHESIS, SELF-TEST & POCKET CARD
% ============================================================================
\section{Comparative synthesis, self-test \& pocket card}
\secsub{Summary decision rules, rapid self-test, and the complete lecture on a single pocket card.}

\housefig{figures/rnn_variants_comparison.png}{Parameter count comparison across recurrent architecture variants.}

\begin{table}[htbp]
  \centering
  \caption{Recurrent Architecture Selection Matrix}
  \resizebox{\linewidth}{!}{%
  \begin{tabular}{lcccc}
    \toprule
    \textbf{Architecture} & \textbf{Gates} & \textbf{States} & \textbf{Recurrent Params (h=256, d=100)} & \textbf{Best Use Case} \\
    \midrule
    Vanilla RNN & 0 & $h_t$ & $65.5\text{K} + 25.6\text{K} = 91.3\text{K}$ & Short sequences ($T < 10$) \\
    BRNN & 0 & $\overrightarrow{h}_t, \overleftarrow{h}_t$ & $182.7\text{K}$ & Full sequence context (NER, Classification) \\
    GRU & 2 ($r, z$) & $h_t$ & $274.2\text{K}$ & Medium-to-long sequences, fast training \\
    LSTM & 3 ($f, i, o$) & $c_t, h_t$ & $365.6\text{K}$ & Complex long dependencies ($T > 100$) \\
    \bottomrule
  \end{tabular}%
  }
\end{table}

\subsection*{Rapid Self-Test}
\begin{enumerate}
  \item Q1: Why does a $1 \times 1$ convolution in a bottleneck block reduce parameters? \\
        \flipanswer{A1: It reduces channel depth before expensive $3 \times 3$ convolutions, saving up to 8x parameters.}
  \item Q2: What is the primary difference between autoregressive models and latent variable models? \\
        \flipanswer{A2: Autoregressive models use a fixed window of past inputs; latent variable models track state history $h_t$.}
  \item Q3: Why does GRU train faster than LSTM? \\
        \flipanswer{A3: GRU has 2 gates instead of 3, merged state $h_t$, and ~25\% fewer parameters.}
\end{enumerate}

\begin{keytake}
\textbf{POCKET CARD: Session 9 on One Page}
\begin{itemize}
  \item \textbf{CNN Patterns:} Spatial reduction + channel expansion, residual shortcuts ($y = F(x)+x$), bottlenecks ($1 \times 1$), multi-scale Inception.
  \item \textbf{Classic CNNs:} AlexNet (ReLU/Dropout, 60M), VGG-16 ($3 \times 3$ convs, 138M), GoogLeNet (Inception, GAP, 5M).
  \item \textbf{Transfer Learning:} Freeze early features for small/similar datasets; fine-tune for large/different datasets.
  \item \textbf{Vanilla RNN:} $h_t = \tanh(W_{hh} h_{t-1} + W_{hx} x_t + b_h)$. Shared weights across time cause vanishing gradients ($\prod W_{hh}^T$).
  \item \textbf{BPTT:} Unrolls sequence across time to sum parameter gradients.
  \item \textbf{GRU:} $r_t$ (reset), $z_t$ (update), $h_t = z_t \odot h_{t-1} + (1-z_t) \odot \tilde{h}_t$. $3(hd + h^2 + h)$ params.
  \item \textbf{LSTM:} $f_t$ (forget), $i_t$ (input), $o_t$ (output), cell state $c_t$. $4(hd + h^2 + h)$ params. Solves vanishing gradient via constant error carousel.
\end{itemize}
\end{keytake}

\section*{Glossary \& symbol cheat-sheet}
\noindent\textbf{Key terms.}
\begin{description}[leftmargin=8.5em,style=nextline,font=\normalfont\bfseries\color{navy}]
  \item[Bottleneck Block] A $1 \times 1 \to 3 \times 3 \to 1 \times 1$ conv sequence reducing channel dimensionality to save parameters.
  \item[Residual Connection] A skip connection adding input directly to output ($y = F(x) + x$) to preserve gradient flow.
  \item[BPTT] Backpropagation Through Time; unrolling an RNN across time steps to sum gradients.
  \item[Vanishing Gradient] Exponential decay of gradient magnitudes when backpropagating through long time steps.
  \item[Update Gate ($z_t$)] GRU gate controlling the balance between past hidden state and new candidate state.
  \item[Forget Gate ($f_t$)] LSTM gate deciding what percentage of past cell state memory to erase.
\end{description}

\medskip
\noindent\textbf{Symbols at a glance.}
\begin{center}\small
\begin{tabular}{@{}ll@{}}
  \toprule
  \textbf{Symbol} & \textbf{Meaning} \\
  \midrule
  $h_t$ & Hidden state vector at sequence time step $t$ \\
  $c_t$ & Cell state vector (long-term memory) in LSTM \\
  $W_{hh}, W_{hx}, W_{ho}$ & Shared weight matrices for hidden-to-hidden, input-to-hidden, hidden-to-output \\
  $r_t, z_t$ & Reset gate and Update gate vectors in GRU \\
  $f_t, i_t, o_t$ & Forget gate, Input gate, and Output gate vectors in LSTM \\
  $\sigma(z)$ & Sigmoid activation function squashing values to $[0, 1]$ \\
  $\tanh(z)$ & Hyperbolic tangent activation squashing values to $[-1, 1]$ \\
  \bottomrule
\end{tabular}
\end{center}

\end{document}
